\section{Boundaries of the decomposition obstruction}\label{bound:section}

The existence of two immersive decompositions has a sharp rank threshold
for tensors with three factors. It also behaves differently from the
Cohen--Macaulay property under a change of format. We establish these
points before discussing the relation with multiplication tensors.
Throughout this section, immersion means injectivity of the addition
differential on the space of unordered decompositions into distinct
nonzero pure summands.

\subsection{The first possible rank}\label{bound:threshold-subsection}

\begin{theorem}\label{bound:small-rank}
Let $A,B,C$ be finite-dimensional complex vector spaces and let
$1\leq r\leq4$. No tensor in $A\otimes B\otimes C$ has two distinct
immersive unordered decompositions into $r$ nonzero pure summands.
Consequently, among three-factor tensors in arbitrary finite dimensions,
rank five is the first rank admitting this decomposition obstruction.
\end{theorem}

\begin{proof}
The assertion for $r=1$ is immediate. Suppose $2\leq r\leq4$ and
\[
 T=\sum_{i=1}^r a_i\otimes b_i\otimes c_i
\]
is an immersive decomposition. Each list of opposite-factor products
is linearly independent. For example, a nonzero relation
$\sum_i\lambda_i b_i\otimes c_i=0$ and a fixed $0\neq z\in A$
would give the nonzero source tangent vector with components
$\lambda_i z\otimes b_i\otimes c_i$, whose image under the addition
differential is zero. The other two lists are treated in the same way.

It follows that the rank of each factor-column matrix equals the
corresponding flattening rank of $T$. Indeed, the first flattening is
the product of the first factor-column matrix and the transpose of
the full-column-rank matrix of opposite products. The latter transpose
is surjective, so the product has the same image as the former matrix.
We may therefore restrict all three factors to the supports of $T$;
all summands lie in these supports, and any tangent kernel after
restriction would also be a tangent kernel in the original format.
Write their dimensions as $a,b,c$.

If $a,b,c\leq r-1$, the source dimension is $r(a+b+c-2)$, which is
strictly larger than the ambient dimension $abc$. To check this for
all positive $a,b,c$ in the stated range, observe that
\[
 \frac{abc}{a+b+c-2}
\]
is nondecreasing in each variable: its derivative with respect to $a$
is $bc(b+c-2)/(a+b+c-2)^2\geq0$. Its maximum is
$(r-1)^3/(3r-5)$, whose values for $r=2,3,4$ are respectively
$1,2,27/7$, all strictly smaller than $r$. Immersion is thus impossible
unless at least one flattening has rank $r$.

Permute the factors so that $c_1,\ldots,c_r$ are independent, and put
$P_i=a_i\otimes b_i$. The $C$-flattening has image
$L=\langle P_1,\ldots,P_r\rangle$, with the $P_i$ independent.
Let $\mathcal R\subset A\otimes B$ be the rank-at-most-one matrix cone.
Immersion implies
\begin{equation}\label{bound:scalar-tangent}
 L\cap T_{P_i}\mathcal R=\C P_i\qquad(1\leq i\leq r).
\end{equation}
Indeed, suppose $\delta P_i\in L\cap T_{P_i}\mathcal R$ is not
proportional to $P_i$, and write $\delta P_i=\sum_j h_jP_j$.
Cancel $\delta P_i\otimes c_i$ by the tangent variations
$-h_jP_j\otimes c_i$ in the $j$th summands. The $i$th component of
the resulting source tangent vector is
$(\delta P_i-h_iP_i)\otimes c_i\neq0$, while its sum is zero,
contradicting immersion.

Any other $r$-term decomposition has independent $C$-factors and
independent matrix factors spanning the same $L$, because the
$C$-flattening has rank $r$. If its matrix-factor directions are the
original $[P_i]$, comparison in this basis fixes the pure summands.
Otherwise it supplies a rank-one matrix $Q\in L$ whose direction is
different from every $[P_i]$.

The expansion of $Q$ in the basis $P_1,\ldots,P_r$, restricted to its
nonzero coefficients, gives a minimally dependent set of $m$ rank-one
matrices, where $3\leq m\leq r+1\leq5$. Its span $L_0$ has dimension
$m-1$, and every coefficient in its unique relation is nonzero.
Let $U,V$ be the spans of their first and second factors, of dimensions
$u,v$. Writing the relation as
\[
 U_{\mathrm{col}}D V_{\mathrm{col}}^{\mathsf t}=0,
\]
where $D$ is an invertible diagonal $m\times m$ matrix, shows that
$u+v\leq m$: the rank-$v$ image of $D V_{\mathrm{col}}^{\mathsf t}$
is contained in the $(m-u)$-dimensional kernel of
$U_{\mathrm{col}}$.

If $u=1$ or $v=1$, the entire $L_0$ lies in the rank-one-cone tangent
at each of its nonzero rank-one points. Since $\dim L_0=m-1\geq2$,
this contradicts~\eqref{bound:scalar-tangent} at an original $P_i$
in the relation. Otherwise $u,v\geq2$ and $u+v\leq5$. Up to exchange,
the only possibilities are $(2,2)$ and $(2,3)$. The affine tangent to
the rank-one cone in $U\otimes V$ has dimension $u+v-1$, so at any
point $P$ in the relation,
\[
 \dim\bigl(L_0\cap T_P\mathcal R\bigr)
 \geq (m-1)+(u+v-1)-uv.
\]
For $(u,v)=(2,2)$, we have $m\geq4$ and this lower bound is
$m-2\geq2$. For $(2,3)$, we have $m\geq5$ and it is $m-3\geq2$.
Taking an original $P_i$ again contradicts
\eqref{bound:scalar-tangent}. Thus the immersive decomposition is
unique. The rank-five example in
Proposition~\ref{prop:witness-immersion} proves that the bound is sharp.
\end{proof}

Theorem~\ref{bound:small-rank} concerns this particular source of
normalization branches. It does not assert that every secant of rank
at most four is Cohen--Macaulay, or that rank five is the least rank
for failure of arithmetic Cohen--Macaulayness by any method.

\subsection{A Cohen--Macaulay adjacent family}\label{bound:adjacent-subsection}

\begin{proposition}\label{bound:adjacent}
For every integer $c\geq1$, the variety $X_5(3,3,c)$ is
Cohen--Macaulay. For $c\leq5$ it is the entire ambient tensor space.
For $c\geq5$, third flattening identifies it, with reduced structure,
with the variety of $9\times c$ matrices of rank at most five.
\end{proposition}

\begin{proof}
The core decomposition of Lemma~\ref{lemma:witness-core} has addition
differential rank $45$ in format $3\times3\times5$, by
Proposition~\ref{prop:witness-immersion}. A nonzero $45\times45$
derivative minor implies that the image closure has dimension at
least $45$. This is the dimension of the ambient space, so
$X_5(3,3,5)$ is that entire space. Projection of the third factor to
$\C^c$, for $1\leq c\leq5$, proves the corresponding assertion there:
a linear map preserves sums of five pure tensors and cannot increase
border rank.

Now let $c\geq5$. Third flattening puts $X_5(3,3,c)$ inside the
irreducible determinantal variety $D$ of $9\times c$ matrices of
rank at most five, whose dimension is $5(9+c-5)=5(c+4)$.
Pad the core decomposition in its third factor. Each of the $c-5$
new third-factor directions contributes the five independent matrices
$P_1,\ldots,P_5$ to the differential image. These new blocks have
distinct third-factor coordinates and are independent of the original
image. Thus the addition differential has rank
\[
 45+5(c-5)=5(c+4).
\]
The secant and $D$ consequently have the same dimension. Since the
secant is a closed irreducible subvariety of the irreducible $D$,
they are equal as reduced varieties. For $c\geq6$ the
Cohen--Macaulay assertion is the generic determinantal case; in this
format it also follows from
\cite[Theorem~2.4]{CatalisanoGeramitaGimigliano2008}.
For $c=5$ the determinantal condition is vacuous.
\end{proof}

Thus, at fixed rank bound five, first-factor dimension three, and
$c\geq5$, increasing the second-factor dimension from three to four
changes a Cohen--Macaulay variety to a non-$S_2$ variety. This statement
does not establish minimality among all formats. In particular, the
non-$S_2$ result is not a consequence of merely placing the same two
decompositions inside an arbitrary larger tensor space.

\subsection{An obstruction on invertible-slice charts}\label{bound:slice-subsection}

For $T\in A\otimes B\otimes C$ and $\alpha\in A^*$, write
$T_\alpha:B^*\to C$ for the contraction with $\alpha$.

\begin{proposition}\label{bound:invertible-slice}
Let $r\geq1$, let $\dim B=\dim C=r$, and suppose that
$T\in A\otimes B\otimes C$ has an invertible slice $T_\alpha$.
If $T$ has a decomposition into $r$ nonzero pure summands, then either
that unordered decomposition is unique, or every such decomposition
lies in a positive-dimensional family and has a nonzero kernel for
the addition differential. In particular, two distinct immersive
$r$-term decompositions cannot occur on this chart. No conciseness
assumption on the $A$ factor is required.
\end{proposition}

\begin{proof}
In an $r$-term decomposition $T=\sum_i a_i\otimes b_i\otimes c_i$,
invertibility of $T_\alpha$ forces every $\alpha(a_i)$ to be nonzero
and both $(b_i)$ and $(c_i)$ to be bases. Absorb the scalar
$\alpha(a_i)$ into $b_i$ so that $\alpha(a_i)=1$.
For each $\beta\in A^*$, the endomorphism
\[
 X_\beta=T_\beta T_\alpha^{-1}:C\longrightarrow C
\]
is diagonal in the basis $(c_i)$, with eigenvalues $\beta(a_i)$.
The joint eigenspaces are
\[
 E_a=\langle c_i:a_i=a\rangle.
\]
Both their dimensions and their labels $\beta\mapsto\beta(a)$ are
intrinsic to $T$. If the $a_i$ are distinct, every joint eigenspace
is a line. These lines determine the $c_i$ directions, and
$T_\alpha$ determines their paired $b_i$ coefficients. The labels
then fix the normalized $a_i$, proving uniqueness of the pure summands.

If $a_1=a_2=a$, replace the first two summands by
\[
 a\otimes(b_1+t b_2)\otimes c_1
 +a\otimes b_2\otimes(c_2-t c_1).
\]
Their sum is unchanged. The two altered factor lists remain bases,
so the summands remain nonzero and distinct. At $t=0$ the two
source tangent components are respectively
$a\otimes b_2\otimes c_1$ and its negative. They are nonzero and
their sum is zero. Passing to the free finite permutation quotient
does not remove this tangent kernel. The curve is nonconstant modulo
permutations and therefore gives a positive-dimensional family.
An intrinsic joint eigenspace of dimension at least two has repeated
labels in every decomposition, so this argument applies to every one.
\end{proof}

This proposition applies only to tensors admitting $r$ actual pure
summands. It makes no assertion about the Cohen--Macaulay property of
the whole invertible-slice locus, including tensors of ordinary rank
larger than $r$.

\subsection{Concise lifts and the Hilbert--Quot comparison}
\label{bound:moduli-subsection}

A tensor is \emph{concise} if each of its flattenings has full rank in
its factor. The core tensor $T_5$ in~\eqref{eq:T-first}
has flattening ranks $(3,3,5)$: the third rank follows from independence
of the five matrices, while the first two follow from the three
coordinate row directions and the independent rows of $M$, with
$\det M=2$. Hence every slice in any factor has matrix rank at most
three. Padding, invertible changes of basis, and factor permutations
do not change this conclusion. In particular, this tensor has no
invertible slice between two five-dimensional factors.

There is also a direct description of its lifts under the concise
secant construction of~\cite{JagiellaJelisiejew2026}. Let $c\sigma_5$
denote the projective concise secant variety in the ambient format
containing $T_5$.

\begin{proposition}\label{bound:smooth-lifts}
The two displayed decompositions of $T_5$ determine distinct smooth
points of $c\sigma_5$ above $[T_5]$. With suitable framings, the concise
tensor at each point is the unit tensor
$\Delta_5=\sum_{i=1}^5e_i\otimes e_i\otimes e_i$.
\end{proposition}

\begin{proof}
Write either decomposition as
$T_5=\sum_i a_i\otimes b_i\otimes c_i$ in $V_1\otimes V_2\otimes V_3$.
Set $W_1=W_2=W_3=\C^5$ and define
\[
 \widetilde T=\Delta_5,\qquad
 \phi_1(e_i)=a_i,\quad \phi_2(e_i)=b_i,\quad \phi_3(e_i)=c_i.
\]
These data satisfy the intermediate conciseness conditions of
\cite[Proposition~3.20]{JagiellaJelisiejew2026}. The tensor
$\widetilde T$ is concise; after applying $\phi_3$, the first two
flattenings have rank five, because the remaining coordinate basis
vectors separate the nonzero summands. After applying $\phi_2$ and
$\phi_3$, the first flattening has rank five because the products
$b_i\otimes c_i$ are independent: the $c_i$ form a basis of their
support. The final restriction is $T_5$.

In the projective framed construction, these data lie in an open
subset of
\begin{equation}\label{bound:framed-product}
 X_5(5,5,5)^{\mathrm{conc}}
 \times\PP\Hom(W_1,V_1)
 \times\Hom(W_2,V_2)\times\Hom(W_3,V_3).
\end{equation}
Here the first restriction map is projectivized. The superscript
$\mathrm{conc}$ denotes the concise locus of the \emph{affine}
secant. Apart from the equations defining that affine secant, the
conditions are open; see
\cite[Proposition~3.20 and Theorem~3.25]{JagiellaJelisiejew2026}.

The unit tensor is smooth in $X_5(5,5,5)$. Indeed, decompositions whose
three factor lists are bases form a dense open of the source, and
their image is exactly the $\GL_5^3$ orbit of $\Delta_5$. This orbit
is dense and locally closed, hence open; an algebraic group orbit
over $\C$ is smooth. Thus the product~\eqref{bound:framed-product}
is smooth at both framed points. The map from the framed concise
secant to $c\sigma_5$ is a principal $\GL_5^3$ bundle, so smoothness
descends to both unframed points.

To see that the two points are distinct, note that $\Delta_5$ has a
unique five-term decomposition: apply
Proposition~\ref{bound:invertible-slice} with the functional taking
value one on every $e_i$; its five normalized first-factor labels are
distinct. An identification of the two framed data up to
$\GL_5^3$ would therefore permute these five pure summands. Its images
would make the two original decompositions equal, possibly after a
common scalar arising from $[\phi_1]$. That scalar is one since both
sums equal the same nonzero $T_5$. The original decompositions are
distinct, so no such identification exists.
\end{proof}

Under the standard basis identifications, $\Delta_5$ is the
multiplication tensor of the semisimple algebra $\C^5$.
Embedded reduced configurations of five distinct points are smooth
points of a Hilbert scheme of points on a smooth variety: their locus
is the free symmetric-group quotient of the smooth configuration
space of five ordered distinct points. Thus the algebra arising from
these lifts supplies smooth Hilbert points. The two rank-one
configurations themselves are different points of the Hilbert scheme
of the Segre product; the map to the tensor secant forgets that
configuration. In the rank-five formats of Theorem~\ref{thm:main}, the local
Cohen--Macaulay failure of their common image does not transfer through
this description to either configuration space.

For clarity, the conjectural consequences discussed in the literature
must be distinguished from this local statement. In
\cite[\S3.6.1, following Conjecture~3.32]{JagiellaJelisiejew2026},
universal arithmetic Cohen--Macaulayness of Segre secants is related,
via smooth equivalences, to Cohen--Macaulayness assertions for surface
Quot schemes and smoothable Hilbert components. The implications
stated there run from the universal secant assertion; its failure
supplies no negation of those assertions.

Haiman's original Conjecture~5.2.1 in~\cite{Haiman2001} asks for
\emph{normality and Cohen--Macaulayness of the isospectral Hilbert
scheme} over the principal component of
$\operatorname{Hilb}^n(\C^d)$, for every $n,d$. The following paragraph
observes that normality and Cohen--Macaulayness of the principal
component follow by taking the symmetric-group quotient. Consequently,
Cohen--Macaulayness of the smoothable Hilbert component is a weaker
assertion than the original conjecture.

There is also a scheme-structure distinction on the Quot side.
Artin--Hochster's conjecture, as stated in
\cite[Conjecture~3.2]{Stark2026}, concerns the commuting \emph{scheme}
cut out by the entries of $XY-YX$ for two square matrices. The
presentation of surface Quot schemes in
\cite[Theorem~3.2]{Stark2026} uses a stable open subset of this
scheme with framing vectors and a principal general-linear-group
quotient. By contrast, the minimal-border-rank locus defining the
concise secant is explicitly reduced
\cite[Definition~3.14(a)]{JagiellaJelisiejew2026}. Agreement of closed
points alone cannot identify these scheme structures.
Propositions~\ref{bound:invertible-slice} and~\ref{bound:smooth-lifts}
therefore give a limitation of the present decomposition obstruction;
they provide neither a proof nor a counterexample to the Hilbert,
isospectral Hilbert, or Quot Cohen--Macaulay assertions.
